% ભાગ: first-order-logic
% પ્રકરણ: models-theories
% વિભાગ: size-of-structures

\documentclass[../../../include/open-logic-section]{subfiles}

\begin{document}

\olfileid{fol}{mat}{siz}

\olsection{\printtoken{P}{structure}નું કદ વ્યક્ત કરવું}

\begin{explain}
સંરચનાઓના કેટલાક ગુણધર્મો ભાષાના અતાર્કિક સંકેતો વાપર્યા વિના પણ
વ્યક્ત કરી શકીએ છીએ. દાખલા તરીકે, એવાં !!{sentence}s છે જે
!!a{structure}માં ત્યારે અને માત્ર ત્યારે સત્ય હોય જ્યારે
!!{structure}ના !!{domain}માં ઓછામાં ઓછા, વધુમાં વધુ, અથવા ચોક્કસ
આપેલી સંખ્યા~$n$ જેટલા !!{element}s હોય.
\end{explain}

\begin{prop}
!!{sentence}
\begin{multline*}
  !A_{\ge n} \ident \lexists[x_1][\lexists[x_2][\dots\lexists[x_n][{}]]]\\
  \begin{aligned}
  (\eq/[x_1][x_2] \land {}
  \eq/[x_1][x_3] \land \eq/[x_1][x_4] \land \dots \land \eq/[x_1][x_n] \land {}\\
\eq/[x_2][x_3] \land \eq/[x_2][x_4] \land \dots \land {} \eq/[x_2][x_n] \land {} \\
\vdots\\
\eq/[x_{n-1}][x_n])
\end{aligned}
\end{multline*}
!!a{structure}~$\Struct M$માં ત્યારે અને માત્ર ત્યારે સત્ય છે જ્યારે
$\Domain M$માં ઓછામાં ઓછા~$n$ !!{element}s હોય. પરિણામે,
$\Sat{M}{\lnot !A_{\ge n+1}}$ ત્યારે અને માત્ર ત્યારે જ્યારે
$\Domain M$માં વધુમાં વધુ~$n$ !!{element}s હોય.

\end{prop}

\begin{prop}
!!{sentence}
\begin{multline*}
!A_{= n} \ident \lexists[x_1][\lexists[x_2][\dots\lexists[x_n][{}]]] \\
\begin{aligned}
  (\eq/[x_1][x_2] \land {}
  \eq/[x_1][x_3] \land \eq/[x_1][x_4] \land \dots \land \eq/[x_1][x_n] \land {}\\
\eq/[x_2][x_3] \land \eq/[x_2][x_4] \land \dots \land {} \eq/[x_2][x_n] \land {} \\
\vdots\\
\eq/[x_{n-1}][x_n] \land {} \\
\lforall[y][(\eq[y][x_1] \lor \dots \lor \eq[y][x_n]]))
\end{aligned}
\end{multline*}
!!a{structure}~$\Struct M$માં ત્યારે અને માત્ર ત્યારે સત્ય છે જ્યારે
$\Domain M$માં ચોક્કસ~$n$ !!{element}s હોય.
\end{prop}

\begin{prop}
!!{structure} નીચેના ગણનું નિદર્શ ત્યારે અને માત્ર ત્યારે છે જ્યારે
તે અનંત છે:
\[
\{!A_{\ge 1}, !A_{\ge 2}, !A_{\ge 3}, \dots \}.
\]
\end{prop}

એવું કોઈ એક શુદ્ધ તાર્કિક વાક્ય નથી જે~$\Struct M$માં ત્યારે અને માત્ર
ત્યારે સત્ય હોય જ્યારે~$\Domain M$ અનંત હોય. જોકે અતાર્કિક
!!{predicate}s ધરાવતાં એવાં !!{sentence}s આપી શકાય છે જેમનાં નિદર્શો
માત્ર અનંત હોય (ભલે દરેક અનંત !!{structure} તેમનું નિદર્શ ન હોય).
સાન્ત સંરચના હોવાનો ગુણધર્મ અને !!{nonenumerable} સંરચના હોવાનો ગુણધર્મ તો
!!{sentence}sના અનંત ગણ વડે પણ વ્યક્ત કરી શકાતા નથી. આ હકીકતો સઘનતા
અને લેવેનહાઇમ--સ્કોલેમ પ્રમેયોમાંથી ફલિત થાય છે.

\end{document}
